\pdfoutput=1
\documentclass[10pt,twocolumn]{article}

\usepackage[margin=0.72in,columnsep=0.26in,top=0.85in,bottom=0.85in]{geometry}
\usepackage[T1]{fontenc}
\IfFileExists{lmodern.sty}{\usepackage{lmodern}}{}  % scalable fonts (Type1 embed; needed for microtype)
\usepackage[utf8]{inputenc}
\usepackage{amsmath,amssymb}
\usepackage{graphicx}
\usepackage{booktabs}
\usepackage{tabularx}
\usepackage{array}
\IfFileExists{microtype.sty}{\usepackage{microtype}}{}
\usepackage[font=small,labelfont=bf,skip=4pt]{caption}
\usepackage{xcolor}
\usepackage{url}
\usepackage{hyperref}
\definecolor{linkblue}{RGB}{25,85,155}
\hypersetup{colorlinks=true,linkcolor=linkblue,citecolor=linkblue,urlcolor=linkblue,
  pdftitle={KVForge: Progressive Parametric Absorption in Small Language Models},
  pdfauthor={Hemant Joshi}}

% auto-width left-aligned wrapping column for tabularx
\newcolumntype{L}{>{\raggedright\arraybackslash}X}

% pandoc helpers (we do not use pandoc's standalone template)
\providecommand{\tightlist}{\setlength{\itemsep}{0pt}\setlength{\parskip}{0pt}}
\providecommand{\passthrough}[1]{#1}

% Unicode in the source -> pdfLaTeX, via inputenc (portable; no newunicodechar).
\DeclareUnicodeCharacter{00A7}{\S}
\DeclareUnicodeCharacter{00D7}{\ensuremath{\times}}
\DeclareUnicodeCharacter{00B7}{\ensuremath{\cdot}}
\DeclareUnicodeCharacter{2192}{\ensuremath{\rightarrow}}
\DeclareUnicodeCharacter{0394}{\ensuremath{\Delta}}
\DeclareUnicodeCharacter{2248}{\ensuremath{\approx}}
\DeclareUnicodeCharacter{00B1}{\ensuremath{\pm}}
\DeclareUnicodeCharacter{2212}{\ensuremath{-}}
\DeclareUnicodeCharacter{2013}{\textendash}
\DeclareUnicodeCharacter{2014}{\textemdash}
\DeclareUnicodeCharacter{2265}{\ensuremath{\geq}}
\DeclareUnicodeCharacter{2264}{\ensuremath{\leq}}
\DeclareUnicodeCharacter{226A}{\ensuremath{\ll}}
\DeclareUnicodeCharacter{2261}{\ensuremath{\equiv}}
\DeclareUnicodeCharacter{2208}{\ensuremath{\in}}
\DeclareUnicodeCharacter{03C7}{\ensuremath{\chi}}
\DeclareUnicodeCharacter{03B1}{\ensuremath{\alpha}}
\DeclareUnicodeCharacter{03B5}{\ensuremath{\varepsilon}}
\DeclareUnicodeCharacter{211D}{\ensuremath{\mathbb{R}}}
\DeclareUnicodeCharacter{221A}{\ensuremath{\surd}}
\DeclareUnicodeCharacter{00B2}{\textsuperscript{2}}
\DeclareUnicodeCharacter{2074}{\textsuperscript{4}}
\DeclareUnicodeCharacter{2075}{\textsuperscript{5}}
\DeclareUnicodeCharacter{2077}{\textsuperscript{7}}
\DeclareUnicodeCharacter{207B}{\textsuperscript{\ensuremath{-}}}
\DeclareUnicodeCharacter{2080}{\textsubscript{0}}
\DeclareUnicodeCharacter{1D40}{\textsuperscript{\textsf{T}}}
\DeclareUnicodeCharacter{00E0}{\`a}

\graphicspath{{figures/}}
\setlength{\emergencystretch}{2em}
\hyphenation{par-a-met-ric re-triev-al}

% --- vertical-spacing & float tuning (two-column) ---
% twocolumn defaults to \flushbottom, which stretches inter-paragraph glue to
% bottom-align columns; \raggedbottom removes those gaps (natural ragged bottoms).
\raggedbottom

% compact section/subsection headings: \large/\normalsize instead of \Large/\large,
% smaller skips -> no stranded headings, long 5.x titles stay on one line.
\makeatletter
\renewcommand\section{\@startsection{section}{1}{\z@}%
  {-2.2ex \@plus -.5ex \@minus -.2ex}{1.0ex \@plus .2ex}%
  {\normalfont\large\bfseries}}
\renewcommand\subsection{\@startsection{subsection}{2}{\z@}%
  {-1.8ex \@plus -.5ex \@minus -.2ex}{0.6ex \@plus .2ex}%
  {\normalfont\normalsize\bfseries}}
% top-align any unavoidable float page instead of vertically centering it
\setlength\@fptop{0pt}
\setlength\@dblfptop{0pt}
\makeatother

% keep floats close to their text (defaults are ~20pt)
\setlength{\textfloatsep}{9pt plus 2pt minus 2pt}
\setlength{\dbltextfloatsep}{9pt plus 2pt minus 2pt}
\setlength{\floatsep}{8pt plus 2pt minus 2pt}
\setlength{\dblfloatsep}{8pt plus 2pt minus 2pt}
\setlength{\intextsep}{9pt plus 2pt minus 2pt}

% float placement fractions (allow fuller top floats, keep some text per page)
\setcounter{topnumber}{3}
\setcounter{bottomnumber}{2}
\setcounter{totalnumber}{5}
\setcounter{dbltopnumber}{3}
\renewcommand{\topfraction}{0.85}
\renewcommand{\bottomfraction}{0.70}
\renewcommand{\textfraction}{0.10}
\renewcommand{\floatpagefraction}{0.75}
\renewcommand{\dbltopfraction}{0.85}
\renewcommand{\dblfloatpagefraction}{0.65}
